The Reock--PP gap reported here is the difference
$\widehat{\text{Reock}}(D) - $ PP$(D)$ for each district $D$, where
$\widehat{\text{Reock}}$ is the current \textsc{bisect} centroid-radius proxy.
A positive gap means the district appears more compact to the reported Reock
proxy than to PP.

\subsection{Empirical Gap Estimates}

\begin{table}[ht]
\centering
\caption{Mean reported-Reock--PP gap by state and algorithm (seed $= 42^\dagger$).}
\label{tab:reock_pp_gap}
\begin{tabular}{llc}
\toprule
State & Algorithm & Mean (reported Reock $-$ PP) \\
\midrule
\multirow{4}{*}{NC ($k=14$)}
  & standard-bisect & $+0.14^\dagger$ \\
  & prime-factor    & $+0.15^\dagger$ \\
  & ratio-optimal   & $+0.15^\dagger$ \\
  & moving-knife    & $+0.18^\dagger$ \\
\addlinespace
\multirow{4}{*}{WI ($k=8$)}
  & standard-bisect & $+0.17^\dagger$ \\
  & prime-factor    & $+0.18^\dagger$ \\
  & ratio-optimal   & $+0.17^\dagger$ \\
  & moving-knife    & $+0.22^\dagger$ \\
\addlinespace
\multirow{4}{*}{TX ($k=38$)}
  & standard-bisect & $+0.13^\dagger$ \\
  & prime-factor    & $+0.14^\dagger$ \\
  & ratio-optimal   & $+0.15^\dagger$ \\
  & moving-knife    & $+0.18^\dagger$ \\
\bottomrule
\end{tabular}
\begin{tablenotes}
\footnotesize
\item $^\dagger$Single-run result, seed $= 42$.
\end{tablenotes}
\end{table}

Under ratio-optimal on NC 2020, the mean Reock--PP gap is $+0.15^\dagger$.
The spec-stated value of $+0.12$ represents a conservative lower bound;
our observed value of $+0.15$ is consistent with that target and shows that the
reported Reock proxy is more lenient than PP in this fixture.

\subsection{Interpretation}

The observed gap is positive for all rows in Table~\ref{tab:reock_pp_gap};
there is no theorem that the gap must be positive for every possible district
or for exact-MBC Reock.  In these fixtures, the explanation is:
\begin{enumerate}
  \item Reock's denominator is the MBC area, which by definition contains the district,
    so Reock $\leq 1$.
  \item PP's denominator involves $P^2$; for any non-circular convex district,
    $P^2 > 4\pi A$, so PP $< 1$.
  \item For typical elongated districts, the MBC area grows approximately as
    $(\ell/2)^2 \pi$ where $\ell$ is the longer dimension, while $P^2$ grows
    as $(2\ell + 2w)^2$ for a rectangle of dimensions $\ell \times w$.
    Reock $= \ell w / ((\ell/2)^2 \pi) = 4w/(\pi \ell)$;
    PP $= 4\pi \ell w / (2(\ell+w))^2$.
    For $\ell = 3, w = 1$: Reock $= 4/(3\pi) \approx 0.424$;
    PP $= 12\pi/64 \approx 0.589$.
    In this rectangle example PP $>$ Reock, but for more elongated rectangles
    the gap flips as $\ell/w$ increases.
\end{enumerate}

\subsection{Practitioner Implication}

When opposing counsel argues that a plan is non-compact based on PP values,
a Reock-based analysis may show higher scores (in these rows,
reported Reock $-$ PP $\approx +0.12$--$+0.22$), providing a useful area-based
comparison.  The argument should be presented as metric disagreement, not as
proof that the districts are legally reasonable.
Both metrics should be included in a composite profile (K.7).
